\documentclass[11pt]{article}
\usepackage[a4paper,margin=1in]{geometry}
\usepackage{graphicx}
\usepackage{listings}
\usepackage{xcolor}
\usepackage{booktabs}
\usepackage{array}
\usepackage{hyperref}
\usepackage{parskip}
\hypersetup{colorlinks=true, linkcolor=black, urlcolor=blue}

\definecolor{codebg}{RGB}{245,245,245}
\lstdefinestyle{cstyle}{
  language=C,
  backgroundcolor=\color{codebg},
  basicstyle=\ttfamily\footnotesize,
  keywordstyle=\color{blue}\bfseries,
  commentstyle=\color{gray}\itshape,
  stringstyle=\color{teal},
  numbers=left,
  numberstyle=\tiny\color{gray},
  stepnumber=1,
  breaklines=true,
  frame=single,
  tabsize=2,
  showstringspaces=false
}
\lstdefinestyle{outstyle}{
  backgroundcolor=\color{codebg},
  basicstyle=\ttfamily\footnotesize,
  breaklines=true,
  frame=single,
  showstringspaces=false
}

\title{\vspace{-2cm}}
\date{}

\begin{document}
\begin{titlepage}
\centering
\vspace*{1cm}
{\large \textbf{Sri Sivasubramaniya Nadar College of Engineering}}\\
{\small (An Autonomous Institution, Affiliated to Anna University, Chennai)}\\
{\small Kalavakkam -- 603 110}\\[0.3cm]
{\small Department of Computer Science \& Engineering}\\[1.5cm]

\rule{\linewidth}{0.5mm}\\[0.4cm]
{\LARGE \textbf{Computer Networks Laboratory}}\\[0.3cm]
{\Large 1513 -- Computer Networks Lab}\\[0.4cm]
\rule{\linewidth}{0.5mm}\\[1.5cm]

{\huge \textbf{Assignment 5}}\\[0.4cm]
{\Large Simulation of Domain Name Resolution using UDP Socket Programming}\\[2cm]

\begin{center}
\begin{tabular}{ll}
\textbf{Name:} & Sundareswaran Rajaguru \\[0.2cm]
\textbf{Register Number:} & 3122247001048 \\[0.2cm]
\textbf{Year / Semester:} & III Year / V Semester \\[0.2cm]
\textbf{Course:} & 1513 Computer Networks Lab \\[0.2cm]
\textbf{Date of Submission:} & 05 August 2026 \\[0.2cm]
\end{tabular}
\end{center}

\vfill
\end{titlepage}

\section*{1. Objective}
To develop a client/server application using UDP sockets that simulates a DNS server and DNS
client. The client queries the server for domain name resolution and the server responds with
the corresponding IP address. The system also simulates client-side caching of recent lookups
and handles three categories of errors -- non-existent domains (NXDOMAIN), invalid query
format, and request timeouts -- in order to analyse the suitability of UDP (overhead, latency,
reliability) for a DNS-like service.

\section*{2. Server Algorithm}
\begin{enumerate}
  \item Maintain a fixed \texttt{DNSRecord} table (an array of \texttt{domain}$\to$\texttt{ip}
        pairs) acting as the DNS database.
  \item Create a UDP socket, bind it to port 8080.
  \item Loop forever: block on \texttt{recvfrom()} to receive a domain-name query and the
        client's address.
  \item If the query is the special domain \texttt{timeout.com}, deliberately \texttt{sleep(5)}
        and send no reply, to simulate a request that never arrives (timeout demonstration).
  \item Else if the query does not contain a \texttt{'.'} character, it is treated as an
        \emph{invalid format} and the reply \texttt{"ERROR: Invalid Format"} is sent.
  \item Otherwise, linearly search the DNS table for the domain:
  \begin{itemize}
    \item If found, reply with the corresponding IP address.
    \item If not found, reply \texttt{"NXDOMAIN"}.
  \end{itemize}
  \item Send the reply back to the querying client using \texttt{sendto()} and loop.
\end{enumerate}

\section*{3. Client Algorithm}
\begin{enumerate}
  \item Maintain a local \texttt{Cache} array (domain $\to$ IP) that starts empty.
  \item Create a UDP socket (no \texttt{connect()} -- connectionless).
  \item Repeatedly prompt the user for a domain name (or \texttt{exit} to quit).
  \item First check the local cache (\texttt{searchCache()}):
  \begin{itemize}
    \item \textbf{Cache hit} -- print the cached IP immediately, without contacting the server.
    \item \textbf{Cache miss} -- \texttt{sendto()} the domain to the server.
  \end{itemize}
  \item On a cache miss, wait for a reply using \texttt{select()} with a 3-second timeout:
  \begin{itemize}
    \item If \texttt{select()} times out, print \texttt{"Request Timed Out!"}.
    \item Otherwise \texttt{recvfrom()} the reply and interpret it: \texttt{NXDOMAIN} $\Rightarrow$
          domain not found; \texttt{ERROR: Invalid Format} $\Rightarrow$ bad query; else the reply
          is a valid IP, which is displayed and inserted into the client cache.
  \end{itemize}
  \item Display the current cache contents after every query and repeat until the user types
        \texttt{exit}.
\end{enumerate}

\section*{4. Server Code}
\subsection*{4.1 server.c}
\lstinputlisting[style=cstyle]{server.c}

\section*{5. Client Code}
\subsection*{5.1 client.c}
\lstinputlisting[style=cstyle]{client.c}

\section*{6. Output / Test Cases}
The server was started in the background and the client was driven through a single session
that exercises every required scenario: a fresh lookup (cache miss), a repeated lookup of the
same domain (cache hit), a domain absent from the DNS table (NXDOMAIN), a malformed query
(invalid format, no dot), and a query for \texttt{timeout.com} (timeout simulation).

\subsection*{6.1 Server Log}
\begin{lstlisting}[style=outstyle]
====================================
     UDP DNS Server Started
     Listening on Port 8080
====================================

Received Query : google.com
Domain Found
Reply Sent : 142.250.183.14

Received Query : notarealdomain.xyz
Domain Not Found
Reply Sent : NXDOMAIN

Received Query : invalidformat
Reply Sent : ERROR: Invalid Format

Received Query : timeout.com
Simulating Timeout...
\end{lstlisting}

\subsection*{6.2 Client Console Output (all test cases)}
\lstinputlisting[style=outstyle]{output.txt}

Note that the second query for \texttt{google.com} produced a \emph{Cache Hit} on the client
and therefore never reached the server -- this is why the server log shows only one
\texttt{google.com} entry despite two client-side lookups.

\section*{7. Analysis: Suitability of UDP for DNS}
\begin{itemize}
  \item \textbf{Communication overhead:} UDP's 8-byte header and connectionless design avoid the
        three-way handshake and connection-teardown cost of TCP, which matters because a DNS
        query/response is typically a single small round trip -- exactly the workload UDP is
        optimised for.
  \item \textbf{Latency:} Since no connection setup is required, a UDP query reaches the server
        and returns a reply in one round trip, minimising latency for the common case, as observed
        in the near-instantaneous replies for \texttt{google.com} and \texttt{notarealdomain.xyz}.
  \item \textbf{Reliability:} UDP does not guarantee delivery, so the application must implement
        its own timeout and retry logic -- demonstrated here using \texttt{select()} with a
        3-second timeout when the server intentionally withholds a reply for \texttt{timeout.com}.
        In production DNS, if a UDP response does not arrive (or exceeds 512 bytes historically),
        resolvers retry, and can fall back to TCP for larger responses.
  \item \textbf{Caching:} Client-side caching, as implemented here, further reduces the
        reliance on network round trips (and therefore on UDP's lack of guarantees) for repeated
        queries, which mirrors how real resolvers use TTL-based caching to cut both latency and
        server load.
\end{itemize}

\section*{8. Learning Outcomes}
\begin{itemize}
  \item Understood why real-world DNS predominantly uses UDP: low overhead and low latency for
        small, idempotent query/response exchanges outweigh the lack of built-in reliability.
  \item Learned to implement application-level timeout handling over UDP using \texttt{select()},
        since the transport itself provides no delivery guarantee.
  \item Implemented and observed the effect of client-side caching in reducing redundant network
        queries.
  \item Practised handling multiple categories of error conditions (NXDOMAIN, invalid format,
        timeout) end-to-end across a client-server boundary.
  \item Reinforced the connectionless \texttt{sendto()}/\texttt{recvfrom()} programming model and
        contrasted it with the connection-oriented TCP model used in earlier assignments.
\end{itemize}

\end{document}
